\section{Related work}
Reversible, information-conserving automata and the second-order construction are classical
\cite{fredkin1982,toffoli1987,toffoli1990}. Rule 37R and its structures were described by Wolfram
\cite{wolfram2002}, and exactly solvable reversible automata with interacting quasiparticles,
notably the reversible Rule 54, are now well understood
\cite{buca2021,klobas2022,prosen2023,wilkinson2022}. That work analyses scattering and
thermalization, not records left behind once the carriers are gone. Particles as information carriers
and domain walls are formalized in computational mechanics
\cite{hanson1992,hanson1997,hordijk2001,pivato2007,rupe2018}. The closest precedent is local
information dynamics, which identifies background domains as the dominant information storage in
irreversible elementary automata \cite{lizier2012}. Our contribution is interventional and exact
rather than statistical, in line with causal information-flow measures \cite{ay2008,lizier2010caus}.
Collision-based computing uses designed collisions that write stationary memory read by later
gliders \cite{adamatzky2012,martinez2006,martinez2011,martinez2012}. The shadows here are
emergent vacuum-phase residues, not designed localizations. Information left in an environment and
read later has also been studied as stigmergy \cite{klyubin2004} and as outsourced memory through
niche construction \cite{lee2022}.
